
\documentclass[a4paper,oneside,12pt]{article}

\usepackage[utf8]{inputenc}    % make čšž work on input
\usepackage[T1]{fontenc}       % make čšž work on output
\usepackage[english]{babel}    % slovenian language and hyphenation
\usepackage[reqno]{amsmath}    % basic math
\usepackage{amssymb,amsthm}    % math symbols and theorem environments
\usepackage{graphicx}          % images
\usepackage{enumerate}
\usepackage[
  paper=a4paper,
  top=2.5cm,
  bottom=2.5cm,
  left=2.5cm,
  right=2.5cm
% textheight=24cm,
]{geometry}  % page geomerty
\usepackage[bookmarks,
colorlinks=true,
linkcolor=black,
anchorcolor=black,
citecolor=black,
filecolor=black,
menucolor=black,
runcolor=black,
urlcolor=black,
pdfencoding=unicode]{hyperref}
\usepackage{booktabs}
\usepackage{multirow}
\usepackage{longtable}
\hypersetup{pdftitle={ROG corpus documentation}, pdfauthor={Darinka Verdonik}} %To piše v pdf viewerju na vrhu! veri neat
\title{ROG annotation layers and tags}
\author{{}\vspace*{\fill}
Darinka Verdonik\textsuperscript{1}, Kaja Dobrovoljc\textsuperscript{2},
 Peter Rupnik\textsuperscript{3}, Nikola Ljubešić\textsuperscript{3},
 \\ Simona Majhenič\textsuperscript{1}, Jaka Čibej\textsuperscript{2}
\\ \vspace{5cm} \\
\textsuperscript{1} University of Maribor, Faculty of Electrical Engineering
and Computer Science \\
\textsuperscript{2} University of Ljubljana, Faculty of Arts \\
\textsuperscript{3} Department of Knowledge Technologies, Jožef Stefan Institute
\\ \vspace{\fill} \\
}


\begin{document}
\pagenumbering{gobble}
\maketitle
\pagebreak
\pagenumbering{arabic}

\section{Introduction}

This document describes annotation layers and tags of the training corpus of
spoken Slovenian ROG \cite{verdonik2024c}.


ROG corpus consists of:
\begin{description}
  \item[ROG-SST,]  which includes selected Gos 2.1 transcriptions with
    manually assigned lemmas and morphosyntactic tags according to the
    MULTEXT-East annotation scheme, as well as annotations according to the
    Universal Dependencies annotation scheme (i.e. part-of-speech categories,
    morphological features and syntactic dependencies).
    ROG-SST is distributed as CONLL-U format\footnote{
      \url{https://universaldependencies.org/format.html}
    } (2014-2024) (.conllu files).
  \item[ROG-Art,] which includes all the annotation layers from the ROG-SST plus
    additional annotation layers for prosodic units, disfluencies and dialogue
    acts. ROG-Art is distributed as EXMARaLDA format (.EXB files)
    \cite{schmidtworner2014}. In addition to the .EXB files, .EXS files and
    Rog-Art.coma file are added for searching through the annotated corpus in the
    EXMARaLDA EXAKT concordancer, .TRS files \cite{barras2001transcriber} are
    added for viewing the transcriptions without annotations in the Transcriber,
    and .textGrid files \cite{boersma19922022} with additional prosodic
    annotations in the Praat (TeG folder) can be found in the ROG-Art folder.
\end{description}

This documentation provides information on annotation layers and tags for both
ROG-SST and ROG-Art subset.

\section{Data sources and file naming convention}

Both ROG-SST and ROG-Art consist of data published in Gos 2.1 corpus
\cite{verdonik2024a}\footnote{\url{
    http://hdl.handle.net/11356/1863
  }}, which represents a merge of the:
\begin{itemize}
  \item Gos 1  \cite{verdonik2013}\footnote{
          \url{http://hdl.handle.net/11356/1863}} and
  \item selected data from the Artur 1 corpus \cite{verdonik2024b}\footnote{
          \url{http://hdl.handle.net/11356/1772}
        }.
\end{itemize}

\begin{samepage}

  Files preserve Gos 2.1 naming, but in order to enhance clarity and provide additional
  context, prefixes were added to the file names to make them more informative:

  \begin{center}
    \begin{tabular}{|l|p{10cm}|}
      \hline
      \textbf{FILE NAME PREFIX} & \textbf{DESCRIPTION}                            \\ \hline
      \textbf{Rog-Go1}          & The files correspond to the original version
      of the Spoken Slovenian Treebank\cite{dobrovoljcandnivre}, which was
      created by selecting data from the GOS1 corpus based on proportional
      criteria. However, 22 files that were expanded and renamed using the
      Rog-Go2 naming convention are excluded here.                                \\ \hline
      \textbf{Rog-Go2}          & Twenty-two manually selected files containing
      speech samples, which were excluded from the Rog-Go1 subset, were expanded
      by adding additional data from the same speech events from the GOS1 corpus. \\ \hline
      \textbf{Rog-Art}          & The speeches are excerpts taken from the Artur
      speech corpus.                                                              \\ \hline
    \end{tabular}
  \end{center}
\end{samepage}

\section{Capitalization, punctuation, segmentation}
\subsection{Capitalization and punctuation}
\begin{tabular}{|p{3cm}|p{6cm}|p{6cm}|}
  \hline
  \textbf{FILE NAME PREFIX}                                                       &
  \textbf{DESCRIPTION OF CAPITALIZATION AND PUNCTUATION in ROG}                   &
  \textbf{History of capitalization and punctuation in the original corpora (GOS 1, GOS 2.1 and/or ARTUR)} \\ \hline
  \textbf{Rog-Go1 and Rog-Go2}                                                    &
  \multirow{2}{*}{
  \begin{tabular}[l]{p{5.5cm}}
      All three subsets
      (Rog-Go1, Rog-Go2 and Rog-Art) follow identical manually checked
      capitalization and punctuation principles: lowercased transcriptions
      (except for proper nouns), and written-like punctuation, which includes
      both utterance-final punctuation, such as full stops and question marks,
      and utterance-medial punctuation, such as commas.       \\ While punctuation in
      Rog-Art remains identical as in the original ARTUR database, it had to be
      manually added to Rog-Go1 and Rog-Go2 data (see right). \\ Transcription
      and punctuation principles are identical in both CONLL-U and EXB versions
      of the ROG data.
    \end{tabular}} &
  Rog-Go1 and Rog-Go2 transcriptions in the original GOS 1 corpus  follow the
  same truecasing convention (capital letters only for proper names), but a more
  limited set of punctuation (?, !, and …). The same principles were preserved,
  when extending GOS 1 to versions GOS 2.0 and GOS 2.1.                                                    \\ \cline{1-1} \cline{3-3}
  \textbf{Rog-Art}                                                                &   &
  Rog-Art transcriptions in the original ARTUR database adhered to standard
  orthographic conventions for capitalization and punctuation. This includes
  capitalizing both proper names and sentence-initial words, and using both
  sentence-medial and sentence-final punctuation. When adding ARTUR data to
  GOS 2.0, punctuation was removed, but then reintroduced in GOS 2.1, together
  with automatic conversion to truecase capitalization.                                                    \\ \hline
\end{tabular}


\subsection{Segmentation}
\begin{minipage}{\linewidth}
  \begin{center}
    \begin{tabular}{|p{3cm}|p{7cm}|p{7cm}|}
      \hline
      \textbf{FILE NAME PREFIX}                                                                     & \textbf{DESCRIPTION OF SEGMENTATION in ROG} &
      \textbf{History of segmentation in the original corpora (GOS 1, GOS 2.1 and/or ARTUR)}                                                        \\ \hline
      \textbf{\begin{tabular}[c]{p{2cm}}Rog-Go1 \\ and Rog-Go2\end{tabular}} &
      Segmentation in Rog-Go1 and Rog-Go2 has been inherited from the original GOS 1
      and 2.1 corpus (see right) and remained unchanged. The sentence IDs are
      identical to those in GOS 1, ensuring long-term traceability and mapping
      between the datasets.                                                                         &
      GOS 1 was manually segmented into utterances, broadly defined as a
      `syntactically, semantically and prosodically completed unit'. The same
      segmentation was also preserved in GOS 2.0 and 2.1.\footnote{``Izjava je
        osnovna enota govora, ki približno ustreza pojmu povedi v pisnem
        jeziku. Določimo jo tako, da je prozodično, semantično in skladenjsko
        zaokrožena
        enota.
        Pri zelo strnjenem govoru smo pozorni, da izjavo označimo na takem mestu, da
        je v signalu mogoče nedvoumno označiti mejo med izjavama, ne da bi pri tem
        odrezali konec prejšnje ali začetek nove izjave.
        Kjer zgornji vodili omogočata različne odločitve, imajo prednost
        krajše enote.'' \cite{ssjgos}
      }                                                                                                                                             \\\hline
      \textbf{Rog-Art}                                                                              &

      The segmentation of Rog-Art data in EXB format  does not feature any explicit segmentation, as all words are on the same level, direct children of <tier> elements.

      In CONLL-U, Rog-Art has been split into utterances (written-like sentences) to align with the segmentation of Rog-Go1 and Rog-Go2 (see above). This was performed by merging the short segments from the original corpus (see right) based on punctuation marks. Specifically, consecutive segments within each utterance were combined into a single segment until a punctuation mark '.', '?', or '...' was encountered at the end of the first segment, ensuring the segmentation better reflected syntactic structures.

      These merged sentence units are traceable through sentence ID's.                              &
      The manual segmentation in the original Artur database, which was also propagated to Gos 2.0 and Gos 2.1, was mostly prosody-based, i.e. introducing a segment split after each pause, and did not necessarily align with syntactic unit boundaries, such as syntactically complete phrases or sentences.\footnote{``Pomembno je, da segmenti niso predolge enote in da je tam, kjer naredimo mejo segmenta, dovolj premora v govoru, da lahko določimo mejo segmenta, ne da odrežemo del predhodnega ali del naslednjega fonema. /…/ Okvirno velja pravilo, da naj bo premor dolg vsaj 0,2 sekunde, lahko pa tudi manj v naslednjih primerih: (1) če bi bili sicer segmenti predolgi, tj. več kot okoli 10 sekund, (2) če se menjajo govorci ali (3) če bi v segment s hkratnim govorom tako zajeli tudi veliko nehkratnega govora. Vedno pa mora biti med segmenti vsaj malo premora, da je meja med besedami vidna tudi na signalu. /…/ Idealno je, če lahko segment sovpada s stavki ali kratkimi povedmi oz. s semantično in skladenjsko kolikor mogoče zaključenimi, smiselnimi enotami, vendar pogosto to ni mogoče. Glavni vodili za meje med segmenti sta zato:
        • kratek premor v govoru in
        • dolžina segmenta, ki naj ne bo predolga, tj. več kot okoli 10 sekund.'' \cite{verdonikbizjak2023}
      }

      \\ \hline
    \end{tabular}
  \end{center}
\end{minipage}


\section{Lemmas}
Lemmas were first assigned automatically using the CLASSLA-Stanza model for
lemmatisation of standard Slovenian 2.0~\cite{tercon2023} and then manually
validated/corrected. The manual corrections followed the lemmatization
guidelines available at: \url{https://wiki.cjvt.si/books/05-lemmatization}.


\section{Universal part-of-speech categories (upos)}

Universal part-of-speech categories (UPOS column in ROG-SST files) follow the
Universal Dependencies annotation scheme, documented in the following guidelines:
\begin{itemize}
  \item general UD guidelines for part-of-speech (in English, online): \url{https://universaldependencies.org/u/pos/index.html}
  \item general UD guidelines for Slovenian (in English, online): \url{https://universaldependencies.org/sl/}
  \item detailed UD guidelines for Slovenian (in Slovene, standalone document), v1.7: \url{https://wiki.cjvt.si/books/07-universal-dependencies/page/oznacevalne-smernice}
\end{itemize}
\begin{samepage}

  Part-of-speech tags used:
  \begin{center}

    \begin{tabular}{|l|l|}
      \hline
      \textbf{TAG} & \textbf{DESCRIPTION}      \\ \hline
      ADJ          & adjective                 \\ \hline
      ADP          & adposition                \\ \hline
      ADV          & adverb                    \\ \hline
      AUX          & auxiliary                 \\ \hline
      CCONJ        & coordinating conjunction  \\ \hline
      DET          & determiner                \\ \hline
      INTJ         & interjection              \\ \hline
      NOUN         & noun                      \\ \hline
      NUM          & numeral                   \\ \hline
      PART         & particle                  \\ \hline
      PRON         & pronoun                   \\ \hline
      PROPN        & proper noun               \\ \hline
      PUNCT        & punctuation               \\ \hline
      SCONJ        & subordinating conjunction \\ \hline
      SYM          & symbol                    \\ \hline
      VERB         & verb                      \\ \hline
      X            & other                     \\ \hline
    \end{tabular}
  \end{center}
\end{samepage}
As described by Dobrovoljc~\cite{dobrovoljc2024extending}, a semi-automatic
procedure was employed to produce UD morphological annotation, by which the
manually validated Multext-East (JOS) morphosyntactic tags (see below) have been
converted to UD parts-of-speech and feature combinations with the
jos2ud\footnote{\url{https://github.com/clarinsi/jos2ud}} conversion script.
Mistakes and inconsistencies identified during subsequent dependency annotation
and data validation have been manually corrected.


\section{MULTEXT-East morphosyntactic tags (xpos)}

Morphosyntactic tags in ROG-Art and ROG-SST were assigned automatically using
the The CLASSLA-Stanza model for morphosyntactic annotation of standard
Slovenian 2.0\footnote{\url{http://hdl.handle.net/11356/1767}} and then manually
validated/corrected. The tags follow the Slovene Multext-East Specifications
(v6)\footnote{\url{https://nl.ijs.si/ME/V6/msd/html/msd-sl.html}}.

\section{Universal morphological features (feats)}
Universal morphological features (FEATS column in ROG-SST files) follow the
Universal Dependencies annotation scheme, documented in the following guidelines:
\begin{itemize}
  \item general UD guidelines for morphological features (in English, online):
        \url{https://universaldependencies.org/u/feat/index.html}
  \item general UD guidelines for Slovenian (in English, online): \url{
          https://universaldependencies.org/sl/ }
  \item detailed UD guidelines for Slovenian (in Slovene, standalone document),
        v1.7: \url{https://wiki.cjvt.si/books/07-universal-dependencies/page/oznacevalne-smernice}
\end{itemize}
\begin{samepage}
  Tags for morphological features used (found as feature-value pairs, e.g. Tense=Pres):
  \begin{center}
    \begin{tabular}{|l|l|l|}
      \hline
      \textbf{Feature} & \textbf{Values}                   & \textbf{Description}        \\ \hline
      Abbr             & Yes                               & abbreviation                \\ \hline
      Animacy          & Anim, Inanim                      & animacy                     \\ \hline
      Aspect           & Imp, Perf                         & aspect                      \\ \hline
      Case             & Nom, Gen, Dat, Acc, Loc, Ins      & case                        \\ \hline
      Definite         & Ind, Def                          & definiteness or state       \\ \hline
      Degree           & Pos, Cmp, Sup                     & degree                      \\ \hline
      Foreign          & Yes                               & foreign word                \\ \hline
      Gender           & Masc, Fem, Neut                   & gender                      \\ \hline
      Gender{[}psor{]} & Masc. Fem, Neut                   & possessor's gender          \\ \hline
      Mood             & Ind, Imp, Cnd                     & mood                        \\ \hline
      Number           & Sing, Dual, Plur                  & number                      \\ \hline
      Number{[}psor{]} & Sing, Dual, Plur                  & possessor's number          \\ \hline
      NumForm          & Word, Digit, Roman                & numeral form                \\ \hline
      NumType          & Card, Ord, Mult, Sets             & numeral type                \\ \hline
      Person           & 1, 2, 3                           & person                      \\ \hline
      Polarity         & Neg, Pos                          & polarity                    \\ \hline
      Poss             & Yes                               & possessive                  \\ \hline
      PronType         & Prs, Int, Rel, Dem, Tot, Neg, Ind & pronominal type             \\ \hline
      Reflex           & Yes                               & reflexive                   \\ \hline
      Tense            & Pres, Fut                         & tense                       \\ \hline
      Variant          & Bound, Short                      & alternative form of word    \\ \hline
      VerbForm         & Fin, Inf, Sup, Part, Conv         & form of verb or deverbative \\ \hline
    \end{tabular}
  \end{center}
\end{samepage}
As described by Dobrovoljc~\cite{dobrovoljc2024extending}, a semi-automatic
procedure was employed to produce UD morphological annotation, by which the
manually validated Multext-East (JOS) morphosyntactic tags (see below) have been
converted to UD parts-of-speech and feature combinations with the
jos2ud\footnote{\url{https://github.com/clarinsi/jos2ud}} conversion script.
Mistakes and inconsistencies identified during subsequent dependency annotation
and data validation have been manually corrected.

\section{Universal syntactic relations (head and deprel)}
The annotation of syntactic relations between words (HEAD and DEPREL columns in
ROG-SST files) follows the Universal Dependencies annotation scheme, documented in the following guidelines:
\begin{itemize}
  \item general UD guidelines for syntactic relations (in English, online):
        \url{https://universaldependencies.org/u/dep/index.html}
  \item general UD guidelines for Slovenian (in English, online):
        \url{https://universaldependencies.org/sl/}
  \item detailed UD guidelines for Slovenian (in Slovene, standalone document),
        v1.7: \url{https://wiki.cjvt.si/books/07-universal-dependencies/page/oznacevalne-smernice }
\end{itemize}
\begin{samepage}
  Tags for dependency relations used (without subtypes of the type cc:preconj):


  \begin{longtable}{|l|l|}
    \hline
    \textbf{Tag} & \textbf{Description}               \\ \hline
    \endhead
    %
    acl          & clausal modifier of noun           \\ \hline
    advcl        & adverbial clause modifier          \\ \hline
    advmod       & adverbial modifier                 \\ \hline
    amod         & adjectival modifier                \\ \hline
    appos        & appositional modifier              \\ \hline
    aux          & auxiliary verb                     \\ \hline
    case         & case marking preposition           \\ \hline
    cc           & coordinating conjunction           \\ \hline
    ccomp        & clausal complement                 \\ \hline
    conj         & conjunct                           \\ \hline
    cop          & copula verb                        \\ \hline
    csubj        & clausal subject                    \\ \hline
    dep          & unspecified dependency             \\ \hline
    det          & determiner                         \\ \hline
    discourse    & discourse element                  \\ \hline
    dislocated   & dislocated element                 \\ \hline
    expl         & expletive                          \\ \hline
    fixed        & fixed multi-word expression        \\ \hline
    flat         & flat multi word-expression         \\ \hline
    goeswith     & disjointed token                   \\ \hline
    iobj         & indirect object                    \\ \hline
    list         & list                               \\ \hline
    mark         & marker (subordinating conjunction) \\ \hline
    nmod         & nominal modifier                   \\ \hline
    nsubj        & nominal subject                    \\ \hline
    nummod       & numeric modifier                   \\ \hline
    obj          & (direct) object                    \\ \hline
    obl          & oblique nominal (adjunct)          \\ \hline
    orphan       & dependent of missing parent        \\ \hline
    parataxis    & parataxis                          \\ \hline
    punct        & punctuation symbol                 \\ \hline
    reparandum   & overridden disfluency              \\ \hline
    root         & root element                       \\ \hline
    vocative     & vocative                           \\ \hline
    xcomp        & open clausal complement            \\ \hline
  \end{longtable}
\end{samepage}
As described by Dobrovoljc~\cite{dobrovoljc2024extending}, a semi-automatic
procedure was employed to produce UD dependency annotation, involving
automatically pre-parsed transcriptions with the customized Slovenian Trankit
model~\cite{krsnikdobrovoljc}, which were then manually checked by 2-3
independent annotators.
\section{Additional}
In this tier corrections of transcription are annotated, e.g. when a word or
filler in the transcription is missing, e.g. `nnn' means that filler `nnn' is
missing in the defined segment.

Furthermore, lengthened phonemes in case of the lengthening are marked with
sign~(:) after the phoneme, e.g. `o(:)d' means that phoneme `o' in the word
`od' is lengthened.
\section{Disfluency annotations}

\bibliographystyle{plain}
\bibliography{references}
\end{document}